\documentclass[10pt,a4paper]{amsart}
\usepackage[margin=19mm]{geometry}
\usepackage{amsmath,amssymb,mathtools}
\usepackage[scr=rsfs,cal=euler]{mathalfa}
\usepackage{xfrac}
\usepackage[unicode,hidelinks,bookmarks=false]{hyperref}
\newcommand{\ie}{i.e.\ }
\newcommand{\cf}{cf.\ }
\newcommand{\resp}{respectively\ }
\newcommand{\Subsection}{\subsection}
\providecommand{\nobreakdash}{\nobreak\hbox{-}}
\setlength{\emergencystretch}{2em}
\multlinegap0pt
\usepackage{bm}
\usepackage{bbm}
\usepackage{dsfont}
\usepackage{xspace}
\usepackage{cancel}
\usepackage{mathtools}
\usepackage{amsthm}
\makeatletter
\@ifundefined{thm}{\newtheorem{thm}{Theorem}}{}
\@ifundefined{prop}{\newtheorem{prop}[thm]{Proposition}}{}
\makeatother
\def\IP{{\mathbf P}}
\def\cc{{\mathbf G}_{\mathbf m}}
\def\red{\text{\rm red}}
\title[Continuous Linear Series]{Continuous Linear Series}
\author{Eduardo Esteves, Antonio Nigro, Pedro Rizzo}
\date{Source: arXiv:2510.11623v1 (CC BY 4.0)}
\begin{document}
\maketitle

\noindent\emph{Source and licence.} Continuous Linear Series. Version arXiv:2510.11623v1 (CC BY 4.0), distributed under the Creative Commons Attribution 4.0 International licence, as recorded in the arXiv licence metadata for this exact version and shown on the abstract page https://arxiv.org/abs/2510.11623v1. Full rights notice: licenses/arxiv-2510.11623-CC-BY-4.0.txt.\par\smallskip

\noindent\textit{Editorial scope.} Counted unit: the Thm whose source label is \texttt{chainmaps} in the article (source0.tex lines 663--664, source label \texttt{chainmaps}), delivered together with the article's own complete original proof of it (source0.tex lines 691--869, one \texttt{proof} environment of 179 lines). No mathematics is written, repaired or re-derived: statement and proof are the article's own text line for line. The proof is detached from the statement in the source: the article states the result at lines 663--664 and prints its proof at lines 691--869, and the proof environment's own header names the statement, so the association is the article's own. Required context retained uncounted: hV (lines 521--535); hVV (lines 567--568). Every definition, display and label of those lines is retained unchanged. Omitted from this package and not redistributed: 1--520, 536--566, 569--662, 665--690, 870--2131. Nothing inside a retained excerpt is omitted or altered: every retained body is delivered exactly as the article's lines are written, so no omission receipt of any kind accompanies this package. Citations: the retained excerpts carry no citation command at all, so no bibliography entry of the article is transcribed and no reference is redistributed. Reference adaptation: none was needed and none was made; every reference inside the delivered unit resolves inside it, so no unresolved reference or citation is produced. Printed slips kept literal: no printed word, symbol, hypothesis or number of the counted statement or of its proof was altered. Layout and class: the article's own class is \texttt{amsart} with packages including inputenc, amsmath, amsfonts, amssymb, anysize, amscd, xy, setspace, amsthm, graphicx, pstricks-add, pst-all, caption; the wrapper is the lane's pinned \texttt{amsart} editorial wrapper at 10pt on A4, which restates the theorem-like environments the retained text needs and adds the article's own macro definitions that the retained text needs (\texttt{\textbackslash IP}, \texttt{\textbackslash cc}, \texttt{\textbackslash red}), with extra wrapper packages (bm, bbm, dsfont, xspace, cancel, mathtools). The delivered numbering is the unit's own. Assets: no retained span contains a figure environment, a graphics-inclusion command, a PDF-page inclusion, a table or a TikZ picture; no image file is copied into this package, the package declares no asset dependency and no image file is redistributed with it. Licence: version arXiv:2510.11623v1 of this article is distributed under the Creative Commons Attribution 4.0 International licence, as recorded in the arXiv licence metadata for this exact version and shown on the abstract page https://arxiv.org/abs/2510.11623v1. A full rights notice is delivered at licenses/arxiv-2510.11623-CC-BY-4.0.txt. Characters: every retained excerpt is pure ASCII, so the delivered engine, XeLaTeX with its default Latin Modern text font, reports no missing character.
\begin{prop}\label{hV} Let $V\in G(k)$ be a nonfixed point. 
Let $V_{\infty,1}:=\iota_1^{-1}(V)$ and 
$V_{0,2}:=\iota_2^{-1}(V)$. Let 
$V_{\infty,2}:=\rho_2(V)$ and $V_{0,1}:=\rho_1(V)$. Then $V_{\infty,1}\subsetneqq V_{0,1}$ and $V_{0,2}\subsetneqq V_{\infty,2}$. Furthermore,
\begin{equation}\label{limlim}
\lim_{x\to 0}x*V=V_{0,1}\oplus V_{0,2}\quad \text{and} \quad
\lim_{x\to \infty}x*V=V_{\infty,1}\oplus V_{\infty,2}.
\end{equation}
Finally, $h_V$ is an isomorphism onto its image $\Gamma$, and
$$
\deg \Gamma=
\dim_k V_{0,1}-\dim_k V_{\infty,1}=
\dim_k V_{\infty,2}-\dim_k V_{0,2}.
$$
\end{prop}

\begin{prop} \label{hVV} Let $V,V'\in G(k)$. Assume that $\rho_1(V')=\iota^{-1}_1(V)$ (resp.~$\rho_2(V')=\iota_2^{-1}(V)$). Then the intersection $h_{V}(\IP^1)\,\cap\,h_{V'}(\IP^1)$ is empty or scheme-theoretically a point, the latter if and only if $\rho_2(V)=\iota_2^{-1}(V')$ (resp.~$\rho_1(V)=\iota_1^{-1}(V')$).
\end{prop}

\begin{thm}\label{chainmaps} For each connected $S\in \text{\rm Hilb}^{\phi,\cc}_{H^r_d(X)}(k)$, the induced map $S\to T$ is a chain map.
\end{thm}

\begin{proof}[Proof of Theorem~\ref{chainmaps}] Since the coefficient of $v$ in $\phi$ is zero, we must have that 
$S_\red\subseteq H^r_d(X,L)$ for some invertible sheaf $L$ on $X$ with bidegree $(d,0)$, where ``red" denotes the reduced induced subscheme. Again, $H^r_d(X,L)$ is an invariant subscheme of $G\times T$, where $G:=\text{Grass}(r+1,W)$ for
$$
W:=W_1\oplus W_2:=\Gamma(Y,L_Y)\oplus\Gamma(Z,L_Z(dP)).
$$
Let $p_1$ and $p_2$ denote the projections of $G\times T$ onto the indicated factors.

Since $S$ is connected, so is $S_{\text{red}}$. Thus, since $S$ has linear multivariate polynomial, $S$ is of pure dimension~1, and hence so is $S_\red$. We will prove a series of claims:

\vskip0.2cm

{\bf First Claim:} {\it $S_{\text{\rm red}}$ has exactly one irreducible component $S_i$ satisfying $p_2(S_i)=T_i$ for each $i=0,\dots,d$, and $p_2$ restricts to an isomorphism $S_i\xrightarrow{\cong}T_i$. The remaining components of $S_{\text{\rm red}}$ are collapsed to the $N_\ell$ under $p_2$.}

\begin{proof}[Proof first claim] 
%For the sake of preciseness,we put $S_{-1}:=\emptyset$ and $S_{d+1}:=\emptyset$. 
Since the coefficient of each $s_i$ in $\phi$ is $1$, 
we have $p_{2*}[S_\red]=[T]$. Thus, 
for each $i=0,\dots,d$ there exists a unique irreducible 
component $S_i$ of $S_{\text{red}}$
mapping birationally onto $T_i$ via $p_2$, the
remaining components being collapsed. Since the $T_i$ are rational, so must be the $S_i$, and the induced maps $S_i\to T_i$ are isomorphisms. 
Since $S$ is invariant, and the action of $\cc$ moves the points on
$T_i^*$ transitively for each $i$, it follows that each remaining
irreducible component of $S_\red$ collapses to a point among
$N_0,\dots,N_{d+1}$. 
\end{proof}

For each $i=0,\dots, d$, let $N_{i,0}$ be the unique point on $S_i$ mapped to $N_i$ and $N_{i,\infty}$ that mapped to $N_{i+1}$. For each $i=0,\dots,d+1$, let $S_{i-\frac{1}{2}}:=(S\cap p_2^{-1}(N_i))_\red$. 

\vskip0.2cm

{\bf Second Claim:} {\it The following statements hold:
\begin{enumerate}
\item The $S_{i-\frac{1}{2}}$ are pairwise disjoint.
\item For each $i=0,\dots,d+1$, we have that $S_{i-\frac{1}{2}}$ intersects $S_j$ if and only if $j\in\{i-1,i\}$, in which case $S_{i-\frac{1}{2}}\cap S_j$ is scheme-theoretically a point.
\item For each $i=1,\dots,d$ we have that $S_{i-\frac{1}{2}}$ is a point only if $N_{i-1,\infty}= N_{i,0}$, in which case that is the point on $S_{i-\frac{1}{2}}$.
\end{enumerate}
}

\begin{proof}[Proof second claim] The first statement is clear, as the images $p_2(S_{i-\frac{1}{2}})$ are pairwise disjoint. The third is clear as well, since, by definition, $N_{i-1,\infty}$ and $N_{i,0}$ are on $S_{i-\frac{1}{2}}$, if defined. As for the second, let  $i\in\{0,\dots,d+1\}$. Since $N_i\in T_j$ if and only if $j\in\{i-1,i\}$, we have that $S_{i-\frac{1}{2}}$ intersects $S_j$ if and only if $j\in\{i-1,i\}$. Furthermore, since each $S_j$ is mapped by $p_2$ isomorphically to $T_j$, and the scheme theoretic image of $S_{i-\frac{1}{2}}$ is $N_i$, we must have that $S_{i-\frac{1}{2}}\cap S_j$ is scheme-theoretically a point for each $j\in\{i-1,i\}$.
\end{proof}

{\bf Third Claim:} {\it $S_{i-\frac{1}{2}}$ is connected for each $i=0,\dots,d+1$.}

\begin{proof}[Proof third claim] Let $i\in\{0,\dots,d+1\}$. We will prove the claim through a series of statements.

\vskip0.1cm

{\bf Statement 1:} {\it For each $i=1,\dots,d$, every connected subcurve $C$ of $\overline{S_\red-\{S_{i-1}\cup S_i\}}$ intersecting both $S_{i-1}$ and $S_i$ is contained in $S_{i-\frac{1}{2}}$.}

Indeed, the image $p_2(C)$ is connected and contains points on $T_{i-1}$ and on $T_i$. It does not contain either $T_{i-1}$ or $T_i$ because only $S_{i-1}$ and $S_i$ cover these components, by First Claim. By the same reason, it does not contain points on $T_0\cup\cdots\cup T_{i-2}$ or $T_{i+1}\cup\cdots\cup T_d$. In particular, $p_2(C)$ contains neither $N_{i-1}$ nor $N_{i+1}$, and hence must be a point on both $T_{i-1}$ and $T_i$, thus equal to $N_i$. So $C\subseteq S_{i-\frac{1}{2}}$. 

\vskip0.1cm

{\bf Statement 2:} {\it $S_{i-\frac{1}{2}}$ contains at most one isolated point, which, if it exists, is equal to $N_{0,0}$ if $i=0$, to $N_{d,\infty}$ if $i=d+1$ and to $N_{i-1,\infty}= N_{i,0}$ if $1\leq i\leq d$.}

Indeed, let $Q$ be an isolated point on $S_{i-\frac{1}{2}}$. Then $Q$ belongs to an irreducible component $C$ of $S_\red$ whose image under $p_2$ contains $N_i$. Also, since $Q$ is isolated in $S_{i-\frac{1}{2}}$, we have that $C\not\subseteq S_{i-\frac{1}{2}}$, that is, $p_2(C)\neq N_i$. But then $C=S_{i-1}$ or $C=S_i$, whence $Q=N_{i-1,\infty}$ or $Q=N_{i,0}$, respectively. The former is possible only if $i>0$ and the latter is possible only if $i\leq d$. Hence, if $i=0$ then $C=S_0$ and $Q=N_{0,0}$, whereas if $i=d+1$ then $C=S_d$ and $Q=N_{d,\infty}$.

Assume now that $1\leq i\leq d$. Either $Q=N_{i-1,\infty}$ or $Q=N_{i,0}$. By contradiction, assume $N_{i-1,\infty}\neq N_{i,0}$. We will show directly that neither $N_{i-1,\infty}$ nor $N_{i,0}$ is isolated on $S_{i-\frac{1}{2}}$. Indeed, since $N_{i-1,\infty}\neq N_{i,0}$, we have $S_{i-1}\cap S_i=\emptyset$. Since $S$ is connected, there is a connected subcurve $C$ of $\overline{S_\red-\{S_{i-1}\cup S_i\}}$ intersecting both $S_{i-1}$ and $S_i$. Statement~1 yields $C\subseteq S_{i-\frac{1}{2}}$. Since $C$ intersects $S_i$, there is an irreducible component $C'$ of $C$ intersecting $S_i$. Since $C'\subseteq S_{i-\frac{1}{2}}$, the intersection is $N_{i,0}$, and hence $N_{i,0}$ is not isolated on $S_{i-\frac{1}{2}}$. Analogously, neither is $N_{i-1,\infty}$. 

\vskip0.1cm

{\bf Statement 3:} {\it Let $C$ be a maximal connected subcurve of $S_\red$ such that $C\subseteq S_{i-\frac{1}{2}}$. Then $C$ contains 
$N_{i-1,\infty}$ if $i\geq 1$ and $N_{i,0}$ if $i\leq d$.}

Indeed, reason by contradiction. Without loss of generality, assume $i\leq d$ and $C$ does not contain $N_{i,0}$. Since $p_2(C)=N_i$, this means that $C$ does not intersect $S_i$. Since $S_\red$ is connected, there is a connected subcurve $C'$ of $\overline{S_\red-\{S_i\cup C\}}$ intersecting both $S_i$ and $C$. Let $C''$ be an irreducible component of $C'$ intersecting $C$. Then $N_i\in p_2(C'')$. Since $C'\neq S_i$, either $p_2(C' )=N_i$ or $C''=S_{i-1}$, if $i\geq 1$. The former contradicts the maximality of $C$. Thus $i\geq 1$ and $C''=S_{i-1}$. Since $p_2(C)=N_i$, we must have that $S_{i-1}$ meets $C$ at $N_{i-1,\infty}$. Now, if $S_{i-1}$  meets $S_i$, it must be at a point over $N_i$, but the unique point on $S_{i-1}$ above $N_i$ is $N_{i-1,\infty}$. Since $N_{i-1,\infty}$ lies on $C$ and $C$ does not meet $S_i$, we must have that $S_{i-1}$ does not meet $S_i$. But since $C'$ meets $S_i$ and contains $S_{i-1}$, there is a connected subcurve $C'''\subseteq\overline{C'-S_{i-1}}$ meeting $S_{i-1}$ and $S_i$. By Statement~1, we must have $p_2(C''')=N_i$. Then $C'''$ does not intersect $C$ by the maximality of $C$. In particular, $C'''$ does not contain $N_{i-1,\infty}$. But $C'''$ intersects $S_{i-1}$, and can only intersect $S_{i-1}$ at $N_{i-1,\infty}$, a contradiction.

\vskip0.1cm

{\bf Statement 4:} {\it If $S_{i-\frac{1}{2}}$ is not a point, then 
$S_{i-\frac{1}{2}}$ has no isolated point.} 

Indeed, since $S_{i-\frac{1}{2}}$ has at most one isolated point by Statement~2, if $S_{i-\frac{1}{2}}$ is not a point then it contains a nonisolated point $Q$. Let $C$ be a maximal connected subcurve of $S_\red$ containing $Q$ such that $C\subseteq S_{i-\frac{1}{2}}$. Then Statement~3 says that $C$ contains $N_{i-1,\infty}$ if $i\geq 1$ and $N_{i,0}$ if $i\leq d$. In particular, neither $N_{i-1,\infty}$ for $i\geq 1$ nor $N_{i,0}$ for $i\leq d$ is isolated in $S_{i-\frac{1}{2}}$. But then $S_{i-\frac{1}{2}}$ contains no isolated point by Statement~2.

\vskip0.1cm

{\bf Statement 5:} {\it $S_{i-\frac{1}{2}}$ is connected.}

If $S_{i-\frac{1}{2}}$ is a point, then it is clearly connected. If not, then $S_{i-\frac{1}{2}}$ has no isolated point by Statement~4. It is thus a disjoint union of connected subcurves of $S_\red$. By Statement~3, each such subcurve contains $N_{i-1,\infty}$ if $i\geq 1$ and $N_{i,0}$ if $i\leq d$. But then all these subcurves contain the same point, and hence there is just one of them. In other words, $S_{i-\frac{1}{2}}$ is connected.
\end{proof}

{\bf Fourth Claim:} {\it There are irreducible components $C_0,\dots,C_m$ of $S_\red$ such that
\begin{enumerate} 
\item $C_0=S_0$ and $C_m=S_d$;
\item $C_{i-1}\cap C_i\neq\emptyset$ for $i=1,\dots,m$. 
\item There is an increasing sequence $i_0,i_1,\dots,i_d$ of integers such that $C_{i_\ell}=S_\ell$ for $\ell=0,\dots,d$;
\item For each $\ell=1,\dots,d$, we have $i_{\ell}=i_{\ell-1}+1$ if $S_{\ell-1}$ meets $S_\ell$; otherwise, $i_{\ell}>i_{\ell-1}+1$ and $C_{i_{\ell-1}+1}\cup\cdots\cup C_{i_{\ell}-1}$ is contained in $S_{\ell-\frac{1}{2}}$.
\end{enumerate}}

\begin{proof}[Proof fourth claim]
Start with $C_0:=S_0$. Each time $C_i=S_{\ell-1}$ for $\ell=1,\dots,d$, we do the following: If $S_{\ell}$ intersects $S_{\ell-1}$, we put $C_{i+1}:=S_\ell$. If not, then $N_{\ell-1,\infty}\neq N_{\ell,0}$, and thus $S_{\ell-\frac{1}{2}}$ is not a point by Second Claim. By Third Claim, $S_{\ell-\frac{1}{2}}$ is connected. By Second Claim, there are irreducible components $C$ and $C'$ of $S_{\ell-\frac{1}{2}}$ such that $C$ intersects $S_{\ell-1}$ and $C'$ intersects $S_\ell$. Put $C_{i+1}:=C$. Since $S_{\ell-\frac{1}{2}}$ is connected, there is a sequence of irreducible components of $S_{\ell-\frac{1}{2}}$, each intersecting the next, starting with $C_{i+1}$ and ending with $C'$. Denote these components sequentially by $C_{i+1},\cdots,C_{i+j}$, where $j$ is thus the number of components in the sequence and $C_{i+j}=C'$. Then put $C_{i+j+1}:=S_{\ell}$.
\end{proof} 

We observe that $S_\red$ may not be the union of the $C_i$. Indeed, 
$S_{\ell-\frac{1}{2}}$ may not be the union $C_{i_{\ell-1}+1}\cup\cdots\cup C_{i_{\ell}-1}$.

\vskip0.2cm

{\bf Fifth Claim:} {\it The $p_1(C_i)$ which are orbit closures are either points or smooth rational curves. The sum of their degrees in $G$ is at least $r+1$, with equality only if:
\begin{enumerate}
    \item all the $p_1(C_i)$ are orbit closures, 
    \item for $i,j\in\{0,\dots,m\}$ with $i<j$ such that $p_1(C_i)$ and $p_1(C_j)$ are curves that intersect, $p_1(C_i)\cap p_1(C_j)$ is scheme-theoretically a point, and $p_1(C_{i+1}),\dots,p_1(C_{j-1})$ are all equal to this point.
\end{enumerate}}

\begin{proof}[Proof fifth claim] For each $\ell=1,2$, let $\iota_\ell\:W_\ell\to W$ be the natural injection and $\rho_\ell\:W\to W_\ell$ the natural surjection. Since $S$ is $\cc$-invariant, each $C_i$ is $\cc$-invariant, whence the $p_1(C_i)$ are $\cc$-invariant. 

For each $i=0,\dots,m$, let $V_i\subseteq W$ be the image under $p_1$ of a general point on $C_i$. Put 
\[
a_i:=\min\{\dim_k\rho_1(V)\,|\,V\in p_1(C_i)\}
\quad\text{and}\quad 
A_i:=\max\{\dim_k\rho_1(V)\,|\,V\in p_1(C_i)\}.
\]
Thus $a_i\leq A_i$. If $V_i$ is fixed by the action of $\cc$, then $V_i=\rho_1(V_i)\oplus\rho_2(V_i)$. The same holds for each $V\in p_1(C_i)$ by continuity, whence $a_i=A_i$. In this case, if $p_1(C_i)$ is an orbit closure then $p_1(C_i)$ is the point $V_i$. If $V_i$ is not fixed, then $p_1(C_i)$ is the closure of the orbit of $V_i$ in $G$, and is thus a smooth rational curve of degree $A_i-a_i>0$ by Proposition~\ref{hV}. Furthermore, the boundary of the orbit of $V_i$ consists of two points on $p_1(C_i)$ where $a_i$ and $A_i$ are attained. 
%Now, whether $V_j$ is fixed or not, $p_1(C_i\cap C_j)$ is mapped to the boundary of the orbit of $V_i$ for $j\neq i$. 

Now, for each $i=1,\dots,m$, since $C_{i-1}\cap C_i\neq\emptyset$, it follows that  
the intervals $[a_{i-1},A_{i-1}]$ and $[a_i,A_i]$ intersect, and thus 
\[
I:=\bigcup_{i=0}^m[a_i,A_i]
\]
is an interval contained in $[0,r+1]$.
Now, $C_0=S_0$ and $C_m=S_d$. Since $p_2(N_{0,0})=N_0$ and $p_2(N_{d,\infty})=N_\infty$, we have that
\[
p_1(N_{0,0})\subseteq\Gamma(Y,L|_Y)\oplus\Gamma(Z,L|_Z(-P))\quad
\text{and}\quad 
p_1(N_{d,\infty})\subseteq\Gamma(Y,L|_Y(-(d+1)P))\oplus\Gamma(Z,L|_Z(dP)).
\]
Since $L|_Z(-P)$ and $L|_Y(-(d+1)P)$ have negative degrees, it follows that $A_0=r+1$ and $a_m=0$. 
Thus $I=[0,r+1]$. Then the sum of the degrees of the $p_1(C_i)$ for $i$ such that $V_i$ is nonfixed, which is $\sum_i(A_i-a_i)$, is at least $r+1$.

Assume the sum is exactly $r+1$. Then, first, if $V_i$ is fixed, $p_1(C_i)$ must be a point, $V_i$. Thus, all the $p_1(C_i)$ are orbit closures. Second, the intervals $[a_i,A_i]$ can overlap only at the boundaries. For each $i=1,\dots,m$, since $[a_{i-1},A_{i-1}]$ and $[a_i,A_i]$ intersect, we must have that either 
$a_i=A_{i-1}$ or $A_i=a_{i-1}$. Since $A_0=r+1$ and $a_m=0$, we must have that $A_i=a_{i-1}$ for each $i=1,\dots,m$. 

Let now $i,j\in\{0,\dots,m\}$ with $i<j$ such that $p_1(C_i)$ and $p_1(C_j)$ are curves that intersect. Since $i<j$, we have $A_j\leq a_i$. Since $p_1(C_i)$ and $p_1(C_j)$ intersect, the intervals $[a_i,A_i]$ and $[a_j,A_j]$ intersect as well, and hence $A_j=a_{j-1}=A_{j-1}=\cdots=A_{i+1}=a_i$. It follows that 
$p_1(C_{i+1}),\dots,p_1(C_{j-1})$ are points, actually, the same point $V\in G$, since $C_{i+1}\cup\cdots\cup C_{j-1}$ is connected.  Furthermore, since $A_j=a_i$, and $p_1(C_i)$ and $p_1(C_j)$ are orbit closures that intersect, it follows from Proposition~\ref{hV} that 
there is a unique point of intersection, which is
\[
\iota_1^{-1}(V_i)\oplus\rho_2(V_i)=\rho_1(V_j)\oplus\iota_2^{-1}(V_j).
\]
In addition, Proposition~\ref{hVV} yields that the intersection is transverse. Since 
$C_i$ intersects $C_{i+1}$ and $C_j$ intersects $C_{j-1}$, we must have that $V$ is the intersection.
\end{proof}

Put $S^0:=C_0\cup\dots\cup C_m$. 

\vskip0.2cm

{\bf Sixth Claim:} {\it The sum of the degrees of the $p_1(C_i)$ in $G$ is $r+1$. Each $C_i$ is rational and the map $C_i\to p_1(C_i)$ is either an isomorphism or constant, the latter only if $i=i_\ell$ for some $\ell$. Furthermore, $S^0=S_\red$ and $S$ is generically reduced. In particular, $S_{-\frac{1}{2}}$ and $S_{d+\frac{1}{2}}$ are points.}

\begin{proof}[Proof sixth claim] Since $S$ has Hilbert polynomial $\phi$, and since the coefficient of $u$ in $\phi$ is $r+1$, the products of the degree of the covering $C_i\to p_1(C_i)$ with the degree of $p_1(C_i)$ sum up to at most $r+1$. On the other hand, the degrees of the $p_1(C_i)$ sum up to at least $r+1$ by Fifth Claim. It thus follows that both sums are $r+1$ and that $C_i\to p_1(C_i)$ is either constant or an isomorphism for each $i=0,\dots,m$. However, since $C_i\subseteq G\times T$, the curve $p_1(C_i)$ can only be a point if $p_2(C_i)$ is not one, that is, if $i=i_\ell$ for some $\ell$, in which case $C_i$ is isomorphic to $T_\ell$ under $p_2$. In any case, $C_i$ is a rational curve. 

Finally, the multivariate Hilbert polynomial of $S^0$ has the same linear part as that of $S$. Since $S^0\subseteq S_\red\subseteq S$, it follows that $S^0=S_\red$ and $S$ is generically reduced. 
\end{proof}

\vskip0.2cm

{\bf Seventh Claim:} {\it For each $i,j\in\{0,\dots,m\}$ with $i<j$, the curve $C_i$ intersects $C_j$ only if $j=i+1$, transversally at a single point if so.}

\begin{proof}[Proof seventh claim] Assume $C_i\cap C_j\neq\emptyset$. 
Since the $S_{\ell-\frac{1}{2}}$ are disjoint, and $S_h\cap S_q\neq\emptyset$ only if $|h-q|\leq 1$, we must have 
$i_{\ell-1}\leq i<j\leq i_{\ell}$ for some $\ell$. Now, since $p_2(S_{\ell-\frac{1}{2}})=N_\ell$ scheme-theoretically, and 
$S_{\ell-1}$ is mapped isomorphically to 
$T_{\ell-1}$, 
%for each $p\in\{i_{\ell-1}+1,\dots,i_{\ell}\}$, 
we have that $C_j$ intersects $C_{i_{\ell-1}}$ only if $j=i_{\ell-1}+1$, transversally at $N_{\ell-1,\infty}$ if so. Thus 
the claim holds if $i=i_{\ell-1}$. By analogy, it holds as well if $j=i_\ell$.

Assume $i_{\ell-1}<i<j<i_\ell$. Then the maps $C_q\to p_1(C_q)$ are isomorphisms for all $q$ with $i\leq q\leq j$ by Sixth Claim.
%$p_1(C_i),p_1(C_{i+1}),\dots,p_1(C_j)$ are subcurves of $G$ by Sixth Claim. 
Since $C_i\cap C_j\neq\emptyset$, it follows from Fifth Claim that $j=i+1$ and $p_1(C_i)\cap p_1(C_j)$ is scheme-theoretically a point. Furthermore, since the maps $C_i\to p_1(C_i)$ and $C_j\to p_1(C_j)$ are isomorphisms by Sixth Claim, also $C_i\cap C_j$ is scheme-theoretically a point.
\end{proof}

{\bf Eighth Claim:} {\it The curve $S$ is a chain of rational curves and $p_1\:S\to T$ is a chain map.}

\begin{proof} It follows from Sixth Claim that the $C_i$ are rational, and thus $S^0$ is a chain of rational curves by Seventh Claim. Then $S^0$ is connected of arithmetic genus $0$, and hence the constant part of its multivariate Hilbert polynomial coincides with that of $S$. So does its linear part by Sixth Claim. Hence $S$ and $S^0$ have the same Hilbert polynomial, and since $S^0\subseteq S$, we must have $S^0=S$. So $S$ is a chain of rational curves. First and Fourth Claims now yield that $p_1\colon S\to T$ is a chain map.
\end{proof}

The proof of Theorem~\ref{chainmaps} is now complete.
\end{proof}

\end{document}
